\documentclass[11pt,a4paper]{letter}
\usepackage[a4paper,margin=1in]{geometry}
\usepackage{hyperref}
\usepackage{lmodern}

\signature{Hoang Anh Nguyen \\ (on behalf of all co-authors)}
\address{Department of Geophysics \\ Colorado School of Mines \\ Golden, CO 80401, USA \\ hoanganh\_nguyen@mines.edu}

\begin{document}

\begin{letter}{Editorial Office \\ Communications AI \& Computing \\ Nature Portfolio}

\opening{Dear Editor,}

We are pleased to submit our manuscript, ``\texttt{quFWI}: Hybrid Quantum-Classical Finite Basis Physics-Informed Neural Networks for Wave Propagation and Full Waveform Inversion,'' for consideration in \emph{Communications AI \& Computing}.

The manuscript presents the first integration of parameterized quantum circuits into domain-decomposed physics-informed neural networks (FBPINNs) for seismic inverse problems. Our main result is that a hybrid quantum-classical velocity network achieves a lower L1 velocity error than a classical FBPINN baseline---and than 15 classical hyperparameter variants---in approximately $8\times$ fewer training iterations, while using $\sim$$33$\% fewer trainable parameters, on a laterally-varying synthetic FWI benchmark. The quantum circuit is implemented as a JAX-native differentiable statevector simulator, enabling end-to-end automatic differentiation through the classical PINN, the quantum circuit, and the physics-informed loss, without the overhead of parameter-shift rules.

We believe this work falls squarely within the scope of \emph{Communications AI \& Computing}, in particular the \emph{Hybrid quantum-classical systems}, \emph{AI for scientific discovery}, and \emph{Computational modeling in physics} areas. It represents a concrete demonstration of how near-term quantum parameterization can improve an important scientific inverse problem---full waveform inversion---with broader applicability to medical ultrasound tomography, photoacoustic imaging, and non-destructive evaluation.

All data, code, and reproduction scripts are openly available at \url{https://github.com/x-repos/quFWI}, following FAIR principles. We frame the convergence advantage as an improvement in optimization iterations on classically-simulated circuits, not as a claim of quantum computational speedup; the manuscript discusses limitations (classical simulation, 2D acoustic setup, single source) and the path to hardware validation.

We confirm that this manuscript has not been published elsewhere and is not under consideration by another journal. We suggest the following reviewers with expertise in physics-informed neural networks, FWI, and quantum machine learning:
% FILL: suggest 3-4 reviewers with name, affiliation, and email.

Thank you for considering our submission.

\closing{Sincerely,}

\end{letter}

\end{document}
